% Compile with XeLaTeX on Overleaf. Authors/affiliations intentionally blank.
\documentclass[manuscript,review,nonacm]{acmart}
\usepackage{xeCJK}
% Optional local font files are used only for reproducible local verification.
% Overleaf provides the same Noto family; no font upload is required there.
\IfFileExists{fonts/NotoSerifCJKkr-Regular.otf}{
  \setCJKmainfont[Path=fonts/,BoldFont=NotoSerifCJKkr-Bold.otf]{NotoSerifCJKkr-Regular.otf}
  \setCJKsansfont[Path=fonts/,BoldFont=NotoSerifCJKkr-Bold.otf]{NotoSerifCJKkr-Regular.otf}
}{
  \setCJKmainfont{Noto Serif CJK KR}
  \setCJKsansfont{Noto Serif CJK KR}
}
\xeCJKsetup{CJKspace=true}
\setcopyright{none}
\settopmatter{printacmref=false,printccs=false,printfolios=true}
\acmConference{}{}{}
\acmDOI{}
\acmISBN{}
\acmYear{2026}
\copyrightyear{2026}
\title{Sling Scroll: 조준과 이동을 분리한 원거리 스크롤 기법}
\author{}
\renewcommand{\shortauthors}{}
\begin{document}
\begin{abstract}
사진 갤러리와 긴 문서처럼 세로 방향으로 탐색하는 데이터에서는 현재 화면에서 멀리 떨어진 목표에 도달하기 위해 스크롤을 반복하거나, 빠른 이동 후 위치를 보정해야 한다. 본 연구는 이러한 반복 조작과 도착점 제어의 부담을 줄이기 위해 조준과 이동을 분리한 터치 기법인 \textbf{Sling Scroll}을 제안한다. 사용자는 콘텐츠 위에서 누른 채 당겨 상대적인 목적지를 지정하고, 화면이 정지한 상태에서 날짜나 페이지를 미리 확인한 뒤 손을 떼어 이동한다. 짧은 당김과 긴 당김에 서로 다른 이동 이득을 주어 정밀 조정과 원거리 이동을 연결한다. 현재 갤러리와 문서 프로토타입을 구현했으며, Flick 및 Scrubber와 비교하는 18명 규모의 사용자 실험을 설계했다. 본 중간 보고서는 구현 및 기능 검증 결과를 제시하고, 주관적 신체 부담과 목표 초과 이동을 검증할 평가 계획 및 설계의 한계를 논의한다. 사용자 실험은 아직 수행하지 않았다.
\end{abstract}
\keywords{스크롤, 터치 상호작용, 원거리 탐색, 목적지 미리보기}
\maketitle

\section{Introduction}

스마트폰 갤러리에서 오래전 사진을 찾거나 긴 문서의 다른 부분으로 이동할 때, 목표가 현재 화면에서 멀수록 이동 동작과 위치 확인을 반복하게 된다. 여기서 문제는 단순히 더 빠르게 스크롤하는 것이 아니라, \textbf{적은 조작으로 멀리 이동하면서도 원하는 곳에 멈추는 것}이다. 본 연구가 다루는 1차원 데이터는 항목의 표시 모양과 무관하게 탐색 위치가 하나의 세로 스크롤 좌표로 결정되는 콘텐츠를 뜻한다. 따라서 사진이 격자로 배치된 갤러리도 연구 범위에 포함된다.

관성 스크롤인 Flick은 익숙한 입력을 제공하지만, 놓는 순간의 속도로 이후 이동이 정해져 목표를 지나치면 추가 보정이 필요하다. 전체 목록 위치에 대응하는 Scrubber는 큰 이동을 지원하지만, 목록이 길수록 손가락의 작은 변위가 더 많은 콘텐츠에 대응한다. 이 때문에 원거리 도달성과 도착점 정밀도를 함께 고려해야 한다.

Sling Scroll은 현재 위치를 기준으로 목적지를 먼저 조준하고 나중에 이동하도록 이 과정을 재구성한다. 사용자는 절대 위치를 정확히 알지 못하더라도 날짜와 같은 단서를 미리보기에서 확인하며 당김량을 조절할 수 있다. 연구 질문은 \textbf{이동 전 목적지 확인이 원거리 탐색의 신체적 부담과 조작 횟수를 줄이면서 목표 초과 이동도 줄일 수 있는가}이다. 이를 위해 본 연구는 상대적 목적지 제어와 미리보기를 결합한 프로토타입을 구현하고, 이동 거리별 비교 평가를 설계하였다.

\section{Related Work}

원거리 탐색의 속도와 제어 문제는 이동 중의 시각적 맥락, 가속 방식, 입력 수단을 바꾸는 방향으로 연구되어 왔다. Igarashi와 Hinckley의 Speed-dependent Automatic Zooming(SDAZ)은 스크롤 속도가 커질수록 화면을 축소해 빠른 이동 중에도 주변 맥락을 제공한다 \cite{igarashi2000sdaz}. SDAZ가 변위로 속도를 제어하고 이동 중 화면 배율을 바꾸는 데 비해, Sling Scroll은 변위로 목적지를 정하고 조준 중에는 콘텐츠 위치와 배율을 유지한다. 따라서 핵심 차이는 속도 제어를 목적지 제어로 바꾸고, 도착점 확인을 이동 이전에 배치한다는 데 있다.

Jeong 등의 Adaptive Kinetic Scrolling(AKS)은 연속 동작의 시간 간격과 남은 데이터량을 이용해 관성 스크롤 속도를 증폭한다 \cite{jeong2018aks}. 기존 플릭을 유지하면서 원거리 이동을 돕지만, 이동 의도를 입력 간격으로 추정한다. Sling Scroll은 사용자가 당김으로 이동량을 조절하고 확정 전에 목적지를 확인하도록 한다. Quinn과 Cockburn의 Zoofing은 압력 입력에 따른 줌과 플릭을 결합해 긴 목록 선택을 지원한다 \cite{quinn2009zoofing}. 이에 비해 Sling Scroll은 압력 입력 없이 일반적인 터치 좌표를 사용하며 화면을 축소하지 않는다.

이러한 차이는 기존 기법보다 우수하다는 결과를 뜻하지 않는다. 본 연구는 \textbf{정지된 화면에서 상대적 목적지를 조준하고, 의미 있는 단서를 확인한 뒤 이동하는 구성}이 반복 조작과 위치 보정의 부담을 어떻게 바꾸는지 검증한다.

\section{Proposed System}

\subsection{조준, 미리보기, 확정}

Sling Scroll의 상호작용은 누르기, 당겨서 조준하기, 놓아서 확정하기로 이루어진다. 콘텐츠 위를 누른 지점이 기준점이 되며 별도의 길게 누르기 대기 시간은 없다. 위로 당기면 아래쪽의 이후 콘텐츠를, 아래로 당기면 위쪽의 이전 콘텐츠를 겨냥한다. 가로 이동은 목적지 계산에 영향을 주지 않는다. 조준하는 동안 콘텐츠는 정지하고, 기준점과 입력 위치를 잇는 탄성 밴드가 당김 상태를 보여 준다.

미리보기 카드는 갤러리에서는 목적지의 날짜와 대표 이미지를, 문서에서는 페이지 번호와 제목을 표시한다. 위치 레일은 현재 위치와 예정된 목적지를 함께 보여 준다. 사용자는 이 단서를 보면서 당김량을 되돌리거나 늘릴 수 있다. 손을 떼면 계산한 위치로 이동하며, 시작점과의 세로 차이가 10 CSS px 이내이면 취소된다. 포인터 입력이 취소되거나 화면 상태가 바뀌는 경우에도 이동을 확정하지 않는다. 따라서 목적지를 탐색하는 동작과 실제 콘텐츠 이동이 분리된다.

\subsection{당김 거리와 이동량의 대응}

짧은 당김의 조절 가능성을 유지하면서 긴 당김의 이동 범위를 늘리기 위해 선형항과 3차항을 결합한다. 현재 위치를 $p$, 최대 스크롤 위치를 $M$, 아래쪽이 양수인 세로 변위를 $\Delta y$라 하면 목적지 $p'$는 다음과 같다.

\[
u=\max(0,|\Delta y|-10),\qquad
p'=\operatorname{clamp}\!\left(p-\operatorname{sgn}(\Delta y)\left[u+M\min(u/L,1)^3\right],0,M\right).
\]

여기서 $L=\max(72,\min(0.46H,250,R))$이며, $H$는 콘텐츠 영역 높이, $R$은 기준점에서 당기는 방향의 화면 경계까지 거리에서 20 px를 뺀 값이다. 길이의 단위는 CSS px이다. 작은 변위에서는 이동량을 세밀하게 바꾸고, 큰 변위에서는 목록의 넓은 범위에 도달하도록 설계했다. 다만 이동 이득이 목록 길이와 시작 위치에도 의존하므로, 그 예측 가능성과 원거리 정밀도는 평가가 필요하다.

\subsection{구현 범위}

HTML, CSS, JavaScript로 20개 날짜 그룹의 이미지 240개와 HTML 문서 80페이지를 구현했다. 이미지는 데모용 SVG이며 날짜는 오래된 순서에서 최신 순서로 배치된다. 현재 버전은 Sling Scroll의 동작을 확인하기 위한 것으로, 일반 터치 스크롤과 휠 입력을 차단하고 짧은 당김을 반복해 위치를 보정한다. 날짜 그룹으로의 자동 정렬, 다중 선택, 비교 기법 및 실험 로그 수집은 구현하지 않았다. 다음 평가는 이 프로토타입에 비교 조건과 과제 실행 기능을 추가한 뒤 수행한다.

\section{Evaluation}

\subsection{연구 가설과 실험 조건}

평가는 조준과 이동의 분리가 실제 탐색 부담과 도착점 제어에 이득을 주는지 확인한다. \textbf{H1}은 Sling Scroll의 기법 블록 단위 주관적 신체 부담이 Flick보다 낮다는 것이며, 거리가 멀어질수록 조작 횟수 감소가 커지는지도 보조적으로 살핀다. \textbf{H2}는 가장 먼 거리에서 Sling Scroll의 실제 목표 초과 이동 횟수가 Scrubber보다 적다는 것이다. 완료 시간을 함께 측정하여 부담 감소와 이동 속도의 상충을 확인한다.

참가자 18명을 대상으로 3(기법) $\times$ 3(거리) 피험자 내 실험을 계획한다. 기법은 기본 드래그·관성 스크롤인 Flick, Flick에 전체 목록 위치와 선형 대응하는 가장자리 핸들을 추가한 Scrubber, 그리고 Sling Scroll이다. Scrubber는 날짜 말풍선을 제공하며, Sling Scroll은 앞서 설명한 미리보기와 짧은 당김으로 위치를 조정한다. 따라서 비교 대상은 각 탐색 방식 전체이며 매핑 함수만의 효과를 분리하지 않는다.

실험용 갤러리는 연속된 100일에 하루 12개씩 총 1,200개 이미지로 확장한다. 거리는 5일, 20일, 80일로 구분하며 실제 화면 수는 동일 기기에서 측정한다. 참가자마다 3기법 $\times$ 3거리 $\times$ 2방향 $\times$ 4반복, 총 72회를 수행한다. 기법 순서 6가지에 참가자를 3명씩 배정하고, 각 블록의 24개 시작-목표 쌍은 동일하게 유지하되 제시 순서를 무작위화한다. 위·아래 방향은 균형을 맞추고 시작점과 목표가 목록 양 끝 5일 이내에 놓이지 않도록 한다.

\subsection{과제, 측정, 분석}

한 시행에서 참가자는 제시된 목표 날짜로 이동해 테두리로 표시된 이미지 한 개를 탭한다. 목표 표시는 시각 탐색 부담을 줄이기 위한 것이며, 날짜 제목과 현재 날짜 표시는 세 기법에 공통으로 제공한다. 시작 버튼 탭부터 올바른 목표 탭까지 완료 시간을 기록하고, 탐색 제스처 수, 손가락 이동 거리, 오탭 및 목표 초과 이동 횟수를 수집한다. 초과 이동은 실제 콘텐츠 위치가 목표 이미지의 가시 범위를 통과해 반대편으로 넘어간 사건으로 정의하며, 조준 중 미리보기만 목표를 넘는 경우는 세지 않는다. 긴 이동에서 중간 프레임이 생략되어도 검출할 수 있도록 이동 전후 좌표를 함께 사용한다.

각 기법 블록 후 NASA-TLX의 신체적 요구(Physical Demand) 점수를 0--100으로 수집하고, 나머지 하위 항목과 마지막 선호 순위·인터뷰를 보조 자료로 활용한다 \cite{nasaTlx}. 이 점수는 주관적 신체 부담이며 장기적 피로나 거리별 부담을 직접 측정하지 않는다. 동일한 스마트폰을 한 손으로 잡고 반대 손 검지로 조작하도록 통일한다. 기법당 연습 6회를 제공하고, 본 실험과 별도인 2명 파일럿에서 난이도와 약 25분의 예상 소요 시간을 점검한다.

H1과 H2는 각각 신체 부담 점수와 80일 조건의 참가자별 평균 초과 이동 횟수에 대한 Wilcoxon 부호순위 검정으로 비교하고, 두 주 검정에는 Holm 보정을 적용한다. 조작 횟수는 기법·거리별 참가자 요약값으로 경향을 분석한다. 완료 시간은 로그 변환 후 가정이 충족되면 반복측정 분산분석을 사용한다. 효과 크기와 신뢰구간을 함께 보고하고, 개별 시행을 독립 참가자로 취급하지 않는다. 18명은 계획 표본이며 특정 크기의 효과를 검출할 수 있다고 단정하지 않는다.

\section{Results}

현재까지의 결과는 제안한 상호작용의 구현과 기능 검증이다. 계산 로직·콘텐츠·QR에 대한 자동 테스트 10개가 모두 통과했으며, 제스처 방향, 취소 구간, 이동량 증가, 목록 경계, 작은 화면 조건, 미리보기 인덱스와 날짜 정렬 등을 점검했다. 이는 입력 이벤트 전체나 실기기 사용성 검증과 구별되며, 테스트 통과를 사용자 성능 향상의 근거로 해석하지 않는다.

사용자 실험은 아직 수행하지 않아 완료 시간, 신체 부담 및 정확도의 비교 결과는 없다. 따라서 H1과 H2의 지지 여부는 판단할 수 없다. 다음 단계는 비교 기법, 1,200개 이미지, 목표 탭 판정 및 로그 수집을 추가하고 파일럿으로 절차를 확정하는 것이다.

\section{Discussion}

Sling Scroll의 설계상 가능성은 이동을 실행하기 전에 목적지를 수정할 수 있다는 데 있다. 조준 중의 오류가 곧바로 콘텐츠 이동으로 이어지지 않으므로, 목표를 지나친 뒤 다시 찾는 부담을 줄일 수 있다. 반면 미리보기를 읽고 당김량을 조정하는 시간이 추가된다. 따라서 조작 횟수나 초과 이동의 감소만으로 우수성을 판단하지 않고, 완료 시간과 정신적 요구를 함께 해석해야 한다.

또한 비선형 매핑은 긴 이동을 지원하는 동시에 큰 당김 구간의 민감도를 높인다. 시작점이 화면 가장자리에 가까우면 충분히 당길 공간이 부족하고, 같은 당김도 목록 길이에 따라 다른 이동량을 만든다. 현재의 임계값과 곡선은 구현을 위한 설계값이며 최적값으로 검증되지 않았다. 후속 연구에서는 매핑과 활성화 위치를 비교하고, 일반 Flick 및 사진 선택과 공존할 수 있는 진입 방식을 검토할 필요가 있다.

계획한 실험은 검지 조작, 균일한 날짜별 이미지 수, 목표가 표시된 갤러리에 한정된다. 따라서 한 손 엄지 사용의 편의성, 촬영 빈도가 불규칙한 실제 사진 모음, 긴 문서의 의미적 탐색으로 결과를 바로 일반화할 수 없다. 특히 현재 비교는 미리보기와 매핑이 결합된 기법 전체를 평가하므로, 효과가 나타나더라도 개별 요소의 기여는 별도 실험으로 분리해야 한다.

\section{Conclusion}

본 연구는 1차원 콘텐츠의 원거리 탐색에서 반복 조작과 도착점 제어를 함께 다루기 위해 Sling Scroll을 제안하였다. 상대적 당김으로 목적지를 조준하고 이동 전에 단서를 확인하는 상호작용을 구현했으며, Flick 및 Scrubber와 비교할 평가를 설계했다. 현재 기여는 프로토타입과 검증 가능한 연구 질문을 구체화한 데 있다. 후속 실험에서는 신체 부담, 목표 초과 이동, 완료 시간을 함께 분석하여 이 설계가 유용한 거리와 조건을 밝히고자 한다.

\bibliographystyle{ACM-Reference-Format}
\bibliography{references}
\end{document}
